% GENERATED FILE. Edit canonical lessons, then run npm run generate:books.
% canonical-source-hash: fnv1a64:f8ce33ddb27d62ed
% canonical-lessons: TE-C380-kritam, TE-C380-antakumundu, TE-C380-appatinunci, TE-C380-mrtadeham, TE-C380-ventadu

\chapter{Ever Since}
\label{ch:te-ever-since}
\hlchaptermodality{380}

\begin{chapteropening}
\textbf{By the end of this chapter:} \emph{I can say ago, before that, from then on, a dead body and to haunt.}

Complete the last lesson of chapter 380: I can say ago, before that, from then on, a dead body and to haunt.
\end{chapteropening}

\section[\texorpdfstring{\te{క్రితం} (kritaṁ)}{క్రితం (kritaṁ)}]{\te{క్రితం} (kritaṁ) — ago}

\label{lesson:TE-C380-kritam}



\begin{quote}
Before the new one: say the Telugu for to petition, then the Telugu for to beg.
\end{quote}

\subsection*{You'll want to know: \te{క్రితం}}
\textbf{\te{క్రితం}} — \emph{kritaṁ} — \textquotedblleft{}ago\textquotedblright{}.

\begin{quote}
\te{పది} \te{సంవత్సరాల} \te{క్రితం} \te{ఇక్కడ} \te{పెద్ద} \te{వరద} \te{వచ్చింది}.
\end{quote}

\begin{quote}
\emph{padi saṁvatsarāla kritaṁ ikkaḍa pedda varada vaccindi.}
\end{quote}

\begin{quote}
Ten years ago a big flood came here.
\end{quote}

It follows the span: \te{రెండు} \te{రోజుల} \te{క్రితం}, two days ago.

\begin{cousinweb}
\textbf{\te{క్రితం}} (\emph{kritaṁ}) is the Telugu word for ago; its origin is not settled.
\end{cousinweb}

\begin{grammarlens}[title={What you've built}]
One more word for telling what happened.
\end{grammarlens}

\subsection*{Your turn}
\begin{itemize}
\raggedright
  \item \emph{Say it:} \emph{kritaṁ}
  \item \emph{Say it:} \emph{kritaṁ}, once more
  \item \emph{Say it:} \emph{padi saṁvatsarāla kritaṁ ikkaḍa pedda varada vaccindi.} — Ten years ago a big flood came here.
  \item \emph{Recall:} say the Telugu for to petition, then the Telugu for to beg, then say \emph{kritaṁ} again
\end{itemize}

\begin{tcolorbox}[breakable,colback=teal!4,colframe=teal!35!black,title={Before you move on}]
What does \te{క్రితం} mean? (\textquotedblleft{}Ago\textquotedblright{}.) Say it once more.
\end{tcolorbox}
\section[\texorpdfstring{\te{అంతకుముందు} (antakumundu)}{అంతకుముందు (antakumundu)}]{\te{అంతకుముందు} (antakumundu) — before that}

\label{lesson:TE-C380-antakumundu}



\begin{quote}
Before the new one: say the Telugu for to beg, then the Telugu for ago.
\end{quote}

\subsection*{You'll want to know: \te{అంతకుముందు}}
\textbf{\te{అంతకుముందు}} — \emph{antakumundu} — \textquotedblleft{}before that\textquotedblright{}.

\begin{quote}
\te{అంతకుముందు} \te{అతను} \te{ఎప్పుడూ} \te{ఇక్కడికి} \te{రాలేదు}.
\end{quote}

\begin{quote}
\emph{antakumundu atanu eppuḍū ikkaḍiki rālēdu.}
\end{quote}

\begin{quote}
Before that he had never come here.
\end{quote}

\te{అంత}, that, -\te{కు}, to, and \te{ముందు}, before.

\begin{cousinweb}
\textbf{\te{అంతకుముందు}} (\emph{antakumundu}) is native Telugu: \emph{anta}, \textquotedblleft{}that much\textquotedblright{} (DEDR 1), with \emph{-ku} and \emph{mundu}, \textquotedblleft{}before\textquotedblright{} (DEDR 5020).
\end{cousinweb}

\begin{grammarlens}[title={What you've built}]
One more word for telling what happened.
\end{grammarlens}

\subsection*{Your turn}
\begin{itemize}
\raggedright
  \item \emph{Say it:} \emph{antakumundu}
  \item \emph{Say it:} \emph{antakumundu}, once more
  \item \emph{Say it:} \emph{antakumundu atanu eppuḍū ikkaḍiki rālēdu.} — Before that he had never come here.
  \item \emph{Recall:} say the Telugu for to beg, then the Telugu for ago, then say \emph{antakumundu} again
\end{itemize}

\begin{tcolorbox}[breakable,colback=teal!4,colframe=teal!35!black,title={Before you move on}]
What does \te{అంతకుముందు} mean? (\textquotedblleft{}Before that\textquotedblright{}.) Say it once more.
\end{tcolorbox}
\section[\texorpdfstring{\te{అప్పటినుంచి} (appaṭinuñci)}{అప్పటినుంచి (appaṭinuñci)}]{\te{అప్పటినుంచి} (appaṭinuñci) — from then on; ever since}

\label{lesson:TE-C380-appatinunci}



\begin{quote}
Before the new one: say the Telugu for ago, then the Telugu for before that.
\end{quote}

\subsection*{You'll want to know: \te{అప్పటినుంచి}}
\textbf{\te{అప్పటినుంచి}} — \emph{appaṭinuñci} — \textquotedblleft{}from then on; ever since\textquotedblright{}.

\begin{quote}
\te{అప్పటినుంచి} \te{ఆ} \te{ఇంట్లో} \te{ఎవరూ} \te{ఉండటం} \te{లేదు}.
\end{quote}

\begin{quote}
\emph{appaṭinuñci ā iṇṭlō evarū uṇḍaṭaṁ lēdu.}
\end{quote}

\begin{quote}
Since then nobody has lived in that house.
\end{quote}

\te{అప్పటి}, of that time, and \te{నుంచి}, from.

\begin{cousinweb}
\textbf{\te{అప్పటినుంచి}} (\emph{appaṭinuñci}) is native Telugu: \emph{appaṭi}, \textquotedblleft{}of that time\textquotedblright{}, from \emph{appuḍu} (DEDR 4559), with \emph{nuñci}, \textquotedblleft{}from\textquotedblright{}.
\end{cousinweb}

\begin{grammarlens}[title={What you've built}]
One more word for telling what happened.
\end{grammarlens}

\subsection*{Your turn}
\begin{itemize}
\raggedright
  \item \emph{Say it:} \emph{appaṭinuñci}
  \item \emph{Say it:} \emph{appaṭinuñci}, once more
  \item \emph{Say it:} \emph{appaṭinuñci ā iṇṭlō evarū uṇḍaṭaṁ lēdu.} — Since then nobody has lived in that house.
  \item \emph{Recall:} say the Telugu for ago, then the Telugu for before that, then say \emph{appaṭinuñci} again
\end{itemize}

\begin{tcolorbox}[breakable,colback=teal!4,colframe=teal!35!black,title={Before you move on}]
What does \te{అప్పటినుంచి} mean? (\textquotedblleft{}From then on; ever since\textquotedblright{}.) Say it once more.
\end{tcolorbox}
\section[\texorpdfstring{\te{మృతదేహం} (mṛtadēhaṁ)}{మృతదేహం (mṛtadēhaṁ)}]{\te{మృతదేహం} (mṛtadēhaṁ) — a dead body}

\label{lesson:TE-C380-mrtadeham}



\begin{quote}
Before the new one: say the Telugu for before that, then the Telugu for from then on.
\end{quote}

\subsection*{You'll want to know: \te{మృతదేహం}}
\textbf{\te{మృతదేహం}} — \emph{mṛtadēhaṁ} — \textquotedblleft{}a dead body\textquotedblright{}.

\begin{quote}
\te{నది} \te{ఒడ్డున} \te{ఒక} \te{మృతదేహం} \te{కనిపించింది}.
\end{quote}

\begin{quote}
\emph{nadi oḍḍuna oka mṛtadēhaṁ kanipiñcindi.}
\end{quote}

\begin{quote}
A dead body was found on the riverbank.
\end{quote}

Plural: \te{మృతదేహాలు}, \emph{mṛtadēhālu}.

\te{మృత}, dead, and \te{దేహం}, a body: a news word.

\begin{cousinweb}
\textbf{\te{మృతదేహం}} (\emph{mṛtadēhaṁ}) is borrowed from Sanskrit: \emph{mṛta}, \textquotedblleft{}dead\textquotedblright{}, and \emph{deha}, \textquotedblleft{}the body\textquotedblright{} (Monier-Williams).
\end{cousinweb}

\begin{grammarlens}[title={What you've built}]
One more word for telling what happened.
\end{grammarlens}

\subsection*{Your turn}
\begin{itemize}
\raggedright
  \item \emph{Say it:} \emph{mṛtadēhaṁ}
  \item \emph{Say it:} \emph{mṛtadēhaṁ}, once more
  \item \emph{Say it:} \emph{nadi oḍḍuna oka mṛtadēhaṁ kanipiñcindi.} — A dead body was found on the riverbank.
  \item \emph{Recall:} say the Telugu for before that, then the Telugu for from then on, then say \emph{mṛtadēhaṁ} again
\end{itemize}

\begin{tcolorbox}[breakable,colback=teal!4,colframe=teal!35!black,title={Before you move on}]
What does \te{మృతదేహం} mean? (\textquotedblleft{}A dead body\textquotedblright{}.) Say it once more.
\end{tcolorbox}
\section[\texorpdfstring{\te{వెంటాడు} (veṇṭāḍu)}{వెంటాడు (veṇṭāḍu)}]{\te{వెంటాడు} (veṇṭāḍu) — to haunt; to hound}

\label{lesson:TE-C380-ventadu}



\begin{quote}
Before the new one: say the Telugu for from then on, then the Telugu for a dead body.
\end{quote}

\subsection*{You'll want to know: \te{వెంటాడు}}
\textbf{\te{వెంటాడు}} — \emph{veṇṭāḍu} — \textquotedblleft{}to haunt; to hound\textquotedblright{}.

\begin{quote}
\te{ఆ} \te{రాత్రి} \te{జ్ఞాపకం} \te{అతన్ని} \te{ఇంకా} \te{వెంటాడుతోంది}.
\end{quote}

\begin{quote}
\emph{ā rātri jñāpakaṁ atanni iṅkā veṇṭāḍutōndi.}
\end{quote}

\begin{quote}
The memory of that night still haunts him.
\end{quote}

Past (I …): \te{వెంటాడాను}, \emph{veṇṭāḍānu}.

Also to hunt or chase: \te{పులి} \te{జింకను} \te{వెంటాడింది}.

\begin{cousinweb}
\textbf{\te{వెంటాడు}} (\emph{veṇṭāḍu}) is native Dravidian: DEDR 5527 lists Telugu \emph{veṇṭa}, \textquotedblleft{}hunting, the chase\textquotedblright{}, and \emph{vēṭ-āḍu}, \textquotedblleft{}to hunt\textquotedblright{}.
\end{cousinweb}

\begin{grammarlens}[title={What you've built}]
One more word for telling what happened.
\end{grammarlens}

\subsection*{Your turn}
\begin{itemize}
\raggedright
  \item \emph{Say it:} \emph{veṇṭāḍu}
  \item \emph{Say it:} \emph{veṇṭāḍu}, once more
  \item \emph{Say it:} \emph{ā rātri jñāpakaṁ atanni iṅkā veṇṭāḍutōndi.} — The memory of that night still haunts him.
  \item \emph{Recall:} say the Telugu for from then on, then the Telugu for a dead body, then say \emph{veṇṭāḍu} again
\end{itemize}

\begin{tcolorbox}[breakable,colback=teal!4,colframe=teal!35!black,title={Before you move on}]
What does \te{వెంటాడు} mean? (\textquotedblleft{}To haunt; to hound\textquotedblright{}.) Say it once more.
\end{tcolorbox}
